\section*{الفصل \ref{ch:b1:oscillators-resonance} --- الهزّازات الميكانيكية: التخميد والرنين}

\begin{solution}{exo:b1:oscillators-resonance:1}
$\omega_0 = \sqrt{80/0.20} = \qty{20}{rad/s}$ و $f_0 = \qty{3.2}{Hz}$ و $T_0 =
\qty{0.31}{s}$؛ و $v_{\max} = A\omega_0 = \qty{1.0}{m/s}$؛ و $E = \tfrac12 kA^2 =
\qty{0.10}{J}$.
\end{solution}

\begin{solution}{exo:b1:oscillators-resonance:2}
$T = 2\pi\sqrt{\ell/g}$: أي \qty{2.0}{s} على الأرض و \qty{4.9}{s} على القمر.
أما دور الكتلة--النابض، $2\pi\sqrt{m/k}$، فلا يتضمن $g$:
فيبقى دون تغيير.
\end{solution}

\begin{solution}{exo:b1:oscillators-resonance:3}
$\tau = 2Q/\omega_0 = 100/(4\pi) = \qty{8.0}{s}$؛ والطاقة $\propto \eu^{-2t/\tau}$
تنتصف عند $t = \tau\ln2/2 = \qty{2.8}{s}$؛ وتبلغ $5\%$ بعد $3\tau = \qty{24}{s}$،
أي نحو $48 \approx Q$ اهتزازة.
\end{solution}

\begin{solution}{exo:b1:oscillators-resonance:4}
$X_0 = F_0/k = \qty{1.0}{cm}$؛ وعند الرنين $\approx QX_0 = \qty{10}{cm}$؛
وعند $3\omega_0$: $X_0/\sqrt{64 + 0.09} = \qty{0.13}{cm}$.
\end{solution}

\begin{solution}{exo:b1:oscillators-resonance:5}
$v = Lf_0 = 8.0 \times 1.3 = \qty{10}{m/s}$ (أي \qty{37}{km/h}). وعند
\qty{60}{km/h}: $f = 16.7/8.0 = \qty{2.1}{Hz}$ و $x_r = 1.6$: ومنه $X/X_g =
\sqrt{(1 + 5.2)/(2.4 + 5.2)} = 0.90$ --- أي عزل ضعيف؛ فالعزل
يقتضي $x_r \gg 1$ (وعند \qty{120}{km/h} يكون $x_r = 3.2$: أي $0.45$).
\end{solution}

\begin{solution}{exo:b1:oscillators-resonance:6}
لا يكون شغل الدفعة موجبًا إلا حين تؤثر في اتجاه الحركة؛
فإذا دُفعت في كل دور تجمّعت الدفعات كلها؛ وبأيّ إيقاع آخر
تتناوب. وفي \omterm{def:b1:transient-regimes:firstorder}{النظام الدائم}: $2\pi E/Q = \qty{25}{J}$ ومنه $E = \qty{80}{J}$؛
و $mg\ell(1 - \cos\theta_{\max}) = 80$: أي $1 - \cos\theta_{\max} = 0.109$،
ومنه $\theta_{\max} = 27^\circ$.
\end{solution}

\begin{solution}{exo:b1:oscillators-resonance:7}
$\delta = \ln(1/0.9) = 0.105$ و $Q = \pi/\delta = 30$؛ و $\Delta f = f_0/Q =
\qty{0.067}{Hz}$.
\end{solution}

\begin{solution}{exo:b1:oscillators-resonance:8}
الانتقال: على الطور ($\varphi = 0$) دون الرنين بكثير، و $-\pi/2$ عند
الرنين، و $-\pi$ (أي معاكس) فوقه بكثير. والسرعة عند الرنين: على
الطور مع القوة. والاستطاعة هي $\vect F\cdot\vect v$: فطور السرعة،
لا طور الانتقال، هو الذي يقرّر مقدار الشغل الذي تبذله
الإقادة.
\end{solution}

\begin{solution}{exo:b1:oscillators-resonance:9}
لدينا $V = \omega X$، وهي تساوي
$\omega(F_0/m)/\sqrt{(\omega_0^2 - \omega^2)^2 + (\omega\omega_0/Q)^2}$؛
واقسم البسط والمقام على $\omega\omega_0/Q$، فينتج عندئذ
$V = (F_0Q/m\omega_0)/\sqrt{1 + Q^2(\omega_0/\omega - \omega/\omega_0)^2}$، مع $F_0Q/m\omega_0 = F_0/\alpha$.
والجذر $\geq 1$، ويساوي $1$ تمامًا عند $\omega = \omega_0$.
\end{solution}

\begin{solution}{exo:b1:oscillators-resonance:10}
$P = F_0^2/2\alpha = \qty{10}{W}$؛ و $V_{\max} = F_0/\alpha = \qty{20}{m/s}$ و
$E = \tfrac12 mV_{\max}^2 = \qty{40}{J}$؛ و $E/P = \qty{4.0}{s}$ و $Q/\omega_0 =
m/\alpha = \qty{4.0}{s}$.
\end{solution}

\begin{solution}{exo:b1:oscillators-resonance:11}
$3\tau = \qty{10}{s}$: ومنه $\tau = \qty{3.3}{s}$ و $Q = \omega_0\tau/2 = 2\pi \times
440 \times 3.3/2 \approx 4600$؛ و $\Delta f = f_0/Q = \qty{0.1}{Hz}$. فهي
تستجيب لتواتر واحد أساسًا وتصدره: أي مرجع دقيق حتى
أجزاء قليلة من $10^4$.
\end{solution}

\begin{solution}{exo:b1:oscillators-resonance:12}
لدينا $m\ddot x = mg - \rho_wgS(h + x)$، أي $-\rho_wgSx$ لأن
$mg = \rho_wgSh$:
ومنه $\omega_0^2 = \rho_wgS/m = g/h$ و $T = 2\pi\sqrt{0.10/9.81} = \qty{0.63}{s}$.
ولا تدخل الكتلة ولا كثافة السائل إلا عبر
عمق التوازن $h$.
\end{solution}

\begin{solution}{pb:b1:oscillators-resonance:1}
\textbf{1.} القوى: الوزن (توازنه الاستطالة السكونية)، والنابض
$-ku$، والمخمّد $-\alpha\dot u$؛ ومنه $m\ddot x = -ku - \alpha\dot u$ مع $\ddot x =
\ddot x_g + \ddot u$.

\textbf{2.} $\ddot u + (\omega_0/Q)\dot u + \omega_0^2u = -\ddot x_g$ مع $\omega_0 =
\sqrt{k/m}$ و $Q = \sqrt{km}/\alpha$.

\textbf{3.} $k = m\omega_0^2 = \pi^2 = \qty{9.9}{N/m}$؛ و $\alpha = m\omega_0/Q =
3.14/0.70 = \qty{4.5}{N.s/m}$.

\textbf{4.} $u = -a_g/\omega_0^2$: فتهبط الكتلة بمقدار ثابت --- أي إن
الآلة مقياس تسارع عند التواتر الصفري (ومقياس ميل).

\textbf{5.} الانتقال النسبي $u(t)$.

\textbf{6.} $-\omega^2\underline U + j\omega(\omega_0/Q)\underline U + \omega_0^2\underline U
= \omega^2X_g$؛ ثم اقسم على $\omega_0^2$.

\textbf{7.} $U/X_g = x_r^2/\sqrt{(1 - x_r^2)^2 + x_r^2/Q^2}$.

\textbf{8.} $U \to X_g$: فالكتلة لا تتحرك في \omterm{thm:b1:newton-dynamics:laws}{المرجع العطالي}
(إذ تبقيها عطالتها ساكنةً بينما يعجز النابض اللين عن جرّها)؛
ويتحرك الهيكل حولها، فيسجّل القلم انتقال الأرض
--- أي مقياس زلازل.

\textbf{9.} $U \approx x_r^2X_g = \omega^2X_g/\omega_0^2 = a_g/\omega_0^2$: أي
مقياس تسارع.

\textbf{10.} $U = QX_g = 0.7X_g$: فلا ذروة؛ والنظامان يلتقيان بسلاسة
ولا ترنّ الآلة.

\textbf{11.} من أجل $Q = 0.7$: مستقيم ميله $2$ يرتفع إلى $1$ قرب
$x_r = 1$، ثم مستوٍ عند $1$؛ ومن أجل $Q = 5$: الشيء نفسه مع ذروة ارتفاعها
$5$ عند $x_r = 1$.

\textbf{12.} $\underline U \to -X_g$: أي $u = -x_g$، فالكتلة ساكنة بينما
يتحرك الهيكل.

\textbf{13.} $\Delta\ell = g/\omega_0^2 = 9.81/(2\pi \times 0.02)^2 = \qty{620}{m}$:
وهو أمر عبثي.

\textbf{14.} $\ell = g/\omega_0^2 = \qty{620}{m}$.

\textbf{15.} بوابة الحديقة: يتأرجح البندول حول محور شبه
شاقولي، فلا يعيده إلا جزء ضئيل من $g$، فيعطي ذراع قصير
دورًا طويلًا. والتغذية الراجعة بالقوة: يبقي وشيعٌ الكتلةَ ثابتة،
ويكون التيار اللازم متناسبًا مع تسارع الأرض --- فالصلابة
إلكترونية وقابلة للضبط.

\textbf{16.} $\Delta\ell = 9.81/\pi^2 = \qty{0.99}{m}$: وهو ممكن.

\textbf{17.} $x_r = 0.2$: ومنه $U/X_g = 0.04/\sqrt{0.92 + 0.08} = 0.040$: أي
\qty{40}{\micro m} --- في نظام مقياس التسارع، ويحتاج إلى تضخيم (برافعة
ضوئية أو بالإلكترونيات).

\textbf{18.} $x_r = 10$: ومنه $100/\sqrt{9801 + 204} \approx 1.00$: أي \qty{1.0}{mm}،
وهو تسجيل أمين للانتقال.

\textbf{19.} $\omega_0 = \qty{6.3e3}{rad/s}$: ومنه $U = g/\omega_0^2 = \qty{2.5e-7}{m}$
--- ويُقرأ تغيّرَ فجوة مكثفة.

\textbf{20.} $U/(a_g/\omega_0^2) = 1/\sqrt{(1 - x_r^2)^2 + x_r^2/Q^2}$؛ ومع
$Q = 0.7$ يتناقص باطّراد ويبقى فوق $0.9$ حتى
$x_r \approx 0.7$: أي نحو \qty{700}{Hz}.

\textbf{21.} $\tau = 2Q/\omega_0 = \qty{0.22}{ms}$: فهي تتابع سقوطًا (يدوم عشرات
الميلي ثانية) دون تأخر.

\textbf{22.} $\tau = 1.4/3.14 = \qty{0.45}{s}$: فبعد قفزة أرضية يستقر
التسجيل في نصف ثانية دون رنين.

\textbf{23.} $Q$ الكبير كان سيجعل الآلة ترنّ عند $f_0$ بعد
كل هزّة (فتسجّل نفسها بدل الأرض) وتبلغ ذروة قرب
$f_0$؛ أما $Q \approx 0.7$ فمستوٍ وسريع.

\textbf{24.} $U = QX_g = 0.7X_g$ مقابل القيمة المثالية $X_g$ (أو $a_g/\omega_0^2
= X_g$): أي أدنى بنسبة $30\%$.

\textbf{25.} قارن تواتر الإشارة بالتواتر $f_0$: فوقه بكثير تكون
الكتلة مرجعًا ثابتًا ويكون $u$ هو الانتقال؛ ودونه بكثير يكون $u$
هو التسارع مقسومًا على $\omega_0^2$.
\end{solution}
